\section{Credo ut intelligam}

Augustine inaugurated the path of confessionalism with his \emph{Credo, ut intelligam}. All of the Axial Age Monotheistic religions developed a ``heart path'' which was mystical. This path was recognized by Gurdjieff as being one of the three paths to transcendence, along with the mind (yoga) and the body (fakir). Although an argument might be made that Eastern Orthodoxy\footnote{\url{http://www.lprww.us/THE\%20SPIRITUAL\%20HEART,\%20GOD\%27S\%20CHANNEL.htm}} kept the theory more pure and the practice more intense, the ``way of the heart'' in all three share many resemblances:

\begin{quotex}
A.M.: When the spiritual heart has awakened, you feel a kind of burning, an influx of energy, a feeling of peace in your chest. When it only starts awakening, you feel pain in the chest. At first it's like heart pain, of-ten quite strong. Those who don't understand what's happening, think it's heart trouble. We've had cases in our commune when people felt acute pain at the time when their spiritual heart began to awaken. Doctors thought it was a heart attack, but cardiograms showed no heart disorders. Objective evidence spoke of good health. When the action of the spiritual heart expands, you may feel pain in your chest or even your shoulders. At first the area of action expands, but then it is localized and forms a kind of sphere in the middle of your chest. After a while the pain goes away and you feel the action only once in a while... 
\end{quotex}


For Islam\footnote{\url{http://www.nuradeen.com/archives/Nuradeen.htm}},

\begin{quotex}
... describes it as an organ in the breast of man which if it is well, the entire being is well; and if it is not, the entire being is not. By the Prophet's indication the \emph{qalb} is not just a physical organ but something like an electromagnetic energy field with its center at the heart. It is a subtler force superimposed on what we call the physical heart. When the heart is turning to a site it will reflect that site. If your heart is connected with a situation, you are likely to reflect that situation. If it is set on something you are likely to attain it, given all the physical possibilities. Wherever your heart turns, it is towards a discernable thing. Its function however, its genetic code is to be ever-turning, never to be fixed upon anything. It is the scanner of objects which gives it the nature of reflectiveness. Its reality, though, is beyond that. The nature of radar and what makes it function is not what it reflects, it is not the image of the aircraft. Its operation is based on wavebands. What makes a heart a heart, a \emph{qalb} a \emph{qalb}, is an entity which is unique, which is \emph{\textbf{Ahad}}, which is not related to anything that the \emph{qalb} reflects and discerns. We can make a holograph of anything, but its origin is not the thing. A heart can reflect anything or any situation, any mental image or thought, but what gives it the power to have that nature of turning, non-attachment and freedom is none of these and it is over and above that, and there is nothing like it. In other words, what is trying to be expressed is that the master of the heart is Allah. The sufferer from the heart's reflection is you and I. And that suffering is part of our endowment. There is no way out of it. It will reflect the good and the bad, the degrees of attachment and non-attachment. If we give it its due we will allow it to be a real \emph{qalb} -- it will turn at all times.
\end{quotex}


Even Evola's group dabbled in ``heart-magic''\footnote{\url{http://cdcruz.wordpress.com/2009/10/08/opus-magicum-fire/}}.

Syncretism would be to indiscriminately blend these together, or to act as if they were all equal. Tradition (to the contrary) looks to what is Truth irregardless of any other consideration except for the wisdom of skillful means and necessary contingencies. Gurdjieff\footnote{\url{http://www.kesdjan.com/exercises/intro.html}} did something like this when he acknowledged the path of the heart as being one of the three primary ways to God, along with mind (yoga) and body (fakir).

In fact, the three stages of Orthodox\footnote{\url{http://www.pelagia.org/htm/b05.en.the\_illness\_and\_cure\_of\_the\_soul.00.htm}} spiritual progress are identical to Evola's:

\begin{enumerate}


\item Purification

\item Illumination

\item Theosis (becoming One with God)
\end{enumerate}


Some traditions (such as the Catholic, and, like Gurdjieff) combined elements of all three major Axial hearth paths into one, with ritual for the body, contemplation for the mind, and prayer or discipline centered in the heart. Bonaventura's\footnote{\url{http://www.franciscan-archive.org/bonaventura/opera/bon05295.html}} \emph{Itinerarium Mentis ad Deum} can be seen as such a permutation, for his times. With Bonaventura, the emphasis is on the mind but blended with a spirituality that is not bookish, but open for any who will make the attempt.

It is likely that Augustinianism, and probably that form as it expresses itself particularly and peculiarly strongly in Christianity, will remain a major attempt by humanity to attain true transcendence. Above and beyond that, it is a developed resource or \emph{upaya} used by Heaven to orient man towards the Good. Although it often becomes little more than ``Platonism for the masses'' and tends to slide towards the Scylla \& Charybdis of syncretism or fundamentalism, a little bit of familiarity with other heart-traditions, or with blended traditions, can begin to assist the mind and body in developing along with the heart, and also open possibilities for improvisation to attain the same end. \emph{We can see that men's experience of the transcendent is one thing, their interpretation often another}. Far from being an inevitable dead-end for fundamentalists or syncretists, the Heart-Path may offer opportunities for new growth and combinations which enhance old patterns enough to make them live again.

The dangers and ugliness potential in such an approach is obvious. The benefits are not quite so plain to see -- yoga remains an attraction to many hip ``seekers'' in the West dissatisfied with spirituality. A heart-path inside a tradition gave some guarantee that seekers would not mistake the Non-being which was real and out of which they were called into existence for God Himself. It gathered energies and potentials for a future consummation.

Mircea Eliade makes the same point in \emph{The Secret of Dr. Honigsberger}, when he has a character remark that asceticism (associated with the heart's ability to reign in the body) is necessary before passing to the other side, not for its own sake, but because the individuality will be exposed without preparation to an entirely unfamiliar world in which it could be lost or destroyed. This was also the basis of much of the advice given in the \emph{Philokalia} about ignoring visions and so forth.

Evola himself was well versed in Catholic dogma, and it seems to have given his mind a cast or mold which served him well both in subtlety, persistence, and also as a defense against the chaotic possibilities which so many encounter as they awaken to the interior life. His awareness of this may well be the basis of his appreciation for what was good in Augustine and the tradition which he founded.


\flrightit{Posted on 2011-04-30 by Logres}

\begin{center}* * *\end{center}

\begin{footnotesize}
\begin{sffamily}
\texttt{Perennial on 2011-04-30 at 01:46 said:}

I personally have my reservations about using St. Augustine a guide, however great a saint he was. St. Augustine had a certain dualistic attitude about the universe, which fits poorly with the holistic model of Tradition. I personally find St. Thomas more consistant and systemic, and therefore more reliable a traditional model, than St. Augustine.

\hfill

\texttt{Logres on 2011-04-30 at 22:11 said:}

Am reliably told that George Heart's \textit{Christianity: Dogmatic Faith \& Gnostic Vivifying Knowledge} has a sympathetic take on Augustine \& his relation to Clement. The only problem with Aquinas is that he seems to take a very abstract, ratiocinated approach to the knowledge of God (as opposed, to say, Bonaventura). A lot of people want to get back to Aquinas, yet somehow the whole thing never gets anywhere. Eager to learn more... 

\hfill

\texttt{Perennial on 2011-05-01 at 03:06 said:}

I believe it is probably more difficult to get back to St. Thomas, more so than St. Augustine, because few of us are taught to think on a level where Thomism would be comfortable discourse. In contrast, I think most people can understand St. Augustine and his pathos, and therefore it makes his way seem more palatable. In seeking a higher path, however, I think St. Thomas would be a better model.

\hfill

\texttt{Cologero on 2011-05-04 at 07:26 said:}

Albert Gleizes makes the claim that Plotinus, Augustine, and Boethius are the intellectual founders of the Middle Ages, a thought which we cannot develop in a comment. (What dualistic attitude?) Yet it is clear that Aquinas is inconceivable without Augustine. Augustine's understanding of the Ideas as existing in the Divine mind and his insight that the Logos described by John is the Logos of the Greeks, recognized interiorly tie him to Tradition.

Augustine saves Aquinas from Aristotle's betrayal of Plato, as Aquinas --- pace Aristotle, and following Augustine --- holds to the ontological priority of ideas over things. Unfortunately, without that, Aquinas leads to nominalism, rationalism and the modern world.

Interestingly, the Thomist Maritain, \emph{en passant}, lists Aristotle, Aquinas and Shankara as his model metaphysicians. However, no one since has been willing or able to integrate Shankara with Thomism. I can only hope that the man who will do that effectively has already been born.

\hfill

\end{sffamily}
\end{footnotesize}

